```latex
\documentclass[12pt]{article}
\usepackage{amsmath,amssymb,amsthm}
\usepackage{hyperref}

\newtheorem{theorem}{Theorem}
\newtheorem{lemma}[theorem]{Lemma}
\newtheorem{proposition}[theorem]{Proposition}
\theoremstyle{definition}
\newtheorem{definition}[theorem]{Definition}

\newcommand{\HFhat}{\widehat{HF}}

\title{Prime summands of Dehn surgeries on L-space knots}
\author{}
\date{}

\begin{document}
\maketitle

\begin{definition}
A knot \(K\subset S^3\) is an \emph{L-space knot} if it admits a positive Dehn surgery
\(S^3_r(K)\) that is an L-space.  A closed oriented \(3\)-manifold \(Y\) is an
\emph{L-space} if \(Y\) is a rational homology sphere and
\[
\dim_{\mathbb F_2}\HFhat(Y;\mathbb F_2)=|H_1(Y;\mathbb Z)|.
\]
\end{definition}

The question asks whether every Dehn surgery on an L-space knot has at most two
prime summands.  The answer is \emph{yes}.  This is a theorem of
Hom--Lidman--Zufelt \cite{HLZ}.

\begin{theorem}[Hom--Lidman--Zufelt \cite{HLZ}]
\label{thm:main}
Let \(K\subset S^3\) be an L-space knot.  Then for every slope \(r\in\mathbb Q\cup\{\infty\}\),
the closed \(3\)-manifold \(S^3_r(K)\) has at most two prime summands.
More precisely, if \(S^3_r(K)\) is reducible, then its prime decomposition has exactly
two summands, one of which is a lens space.
\end{theorem}

The proof of Theorem~\ref{thm:main} rests on the following three established
results.

\begin{lemma}[Gordon--Luecke \cite{GL}]
\label{lem:GL}
Let \(K\subset S^3\) be a knot.  If \(S^3_r(K)\) is reducible, then the surgery
slope \(r\) is an integer.  Moreover, \(S^3_r(K)\) has a lens space summand.
\end{lemma}

\begin{lemma}[Howie--Sayari--Valdez S\'anchez \cite{HSV}]
\label{lem:three}
Let \(K\subset S^3\) be a knot, and suppose that \(S^3_n(K)\) has more than two
prime summands.  Then
\[
S^3_n(K)\cong L(p,a)\#L(q,b)\#Y,
\]
where \(p,q>1\) are coprime, \(L(p,a)\) and \(L(q,b)\) are lens spaces, and
\(Y\) is an integer homology sphere that is not a homotopy sphere.  Furthermore,
the surgery slope satisfies \(|n|=pq\).
\end{lemma}

The next result is the key Floer-theoretic input.  It combines the slope
dichotomy for reducible surgeries on L-space knots with the exclusion of
three-summand surgeries at the exceptional slope.

\begin{theorem}[Hom--Lidman--Zufelt \cite{HLZ}]
\label{thm:HLZ-dichotomy}
Let \(K\) be an L-space knot of genus \(g\).
\begin{enumerate}
\item[(i)] If \(S^3_n(K)\) is reducible for some integer \(n>0\), then either
\(n=2g-1\), or \(K\) is a cable knot and \(n\) is its cabling slope.
\item[(ii)] \(S^3_{2g-1}(K)\) has at most two prime summands.
\end{enumerate}
\end{theorem}

Part (i) is the L-space knot case of the reduction obtained by combining the
L-space surgery slope bound with the Matignon--Sayari bound for reducible
surgeries; part (ii) rules out the remaining potential three-summand surgery at
the slope \(2g-1\).  Both are proved in \cite{HLZ}.

We also use the standard surgery formula for cable knots.

\begin{lemma}[Cable surgery formula]
\label{lem:cable}
Let \(K=C_{p,q}(J)\) be the \((p,q)\)-cable of a knot \(J\), with \(p,q\)
coprime and \(p>1\).  Then the cabling slope is \(pq\), and
\[
S^3_{pq}(K)\cong L(p,q)\# S^3_{p/q}(J).
\]
\end{lemma}

\begin{proof}[Proof of Theorem~\ref{thm:main}]
Let \(M=S^3_r(K)\).

If \(M\) is irreducible, then its prime decomposition has one summand, and there
is nothing to prove.  Suppose therefore that \(M\) is reducible.  By
Lemma~\ref{lem:GL}, the slope \(r\) is an integer; after possibly replacing
\(K\) by its mirror image we may assume \(r=n>0\).  Also by
Lemma~\ref{lem:GL}, \(M\) contains a lens space summand.

We claim that \(M\) cannot have three or more prime summands.  Suppose otherwise.
Then, by Lemma~\ref{lem:three},
\[
M\cong L(p,a)\#L(q,b)\#Y,
\]
with \(p,q>1\) coprime, \(Y\) a nontrivial integer homology sphere, and
\(n=pq\).

On the other hand, since \(K\) is an L-space knot,
Theorem~\ref{thm:HLZ-dichotomy}(i) applies.  Thus either
\(n=2g(K)-1\), or \(K\) is a cable knot and \(n\) is its cabling slope.
The second possibility cannot produce three summands.  Indeed, if
\(K=C_{p,q}(J)\) and \(n=pq\), then by Lemma~\ref{lem:cable},
\[
S^3_n(K)\cong L(p,q)\#S^3_{p/q}(J).
\]
For a nontrivial L-space cable one has \(q>1\), so \(p/q\) is not an integer.
Consequently \(S^3_{p/q}(J)\) is irreducible by the integrality part of
Lemma~\ref{lem:GL}.  Hence in the cabling slope case the surgery has exactly two
prime summands, not three.

Therefore a three-summand surgery on an L-space knot would have to occur at the
slope \(n=2g(K)-1\).  But this is impossible by
Theorem~\ref{thm:HLZ-dichotomy}(ii).  Thus \(M\) has at most two prime
summands.

Finally, if \(M\) is reducible, Lemma~\ref{lem:GL} shows that one of the
summands is a lens space; since there cannot be three or more summands, the
prime decomposition has exactly two summands.  This proves the theorem.
\end{proof}

\begin{thebibliography}{9}

\bibitem{GL}
C. McA. Gordon and J. Luecke,
\emph{Knot surgery and primeness},
Math. Proc. Cambridge Philos. Soc. \textbf{102} (1987), 97--101.
\newline
\href{https://www.cambridge.org/core/journals/mathematical-proceedings-of-the-cambridge-philosophical-society/article/abs/knot-surgery-and-primeness/29D04C8E36EF794942FA2F0925BE1775}{cambridge.org}

\bibitem{HSV}
J. Howie, S. Sayari, and L. Valdez S\'anchez,
\emph{Can Dehn surgery yield three connected summands?},
\newline
\href{https://ar5iv.labs.arxiv.org/html/0908.2697}{ar5iv.labs.arxiv.org}

\bibitem{HLZ}
J. Hom, T. Lidman, and N. Zufelt,
\emph{L-space surgeries, genus bounds, and the cabling conjecture},
\newline
\href{https://exa.ai/library/publication/zl1m8jt741b}{exa.ai}

\end{thebibliography}

\end{document}
```